%%!TEX root = ./UserManual.tex
\chapter{Concepts and Algorithms}
\label{chap:code}


\section{Introduction}
\label{section:ConceptsIntro}

The model population consists of households, schools, workplaces and
communities, which represent a group of people we define as a ``ContactPool''.
Social contacts can only happen within a ContactPool.
When school or work is off, people stay at home and in their primary
community and can have social contacts with the other members.
During other days, people are present in their household, secondary community
and a possible workplace or school.  

We use a $Simulator$ class to organize the activities from the people in the population.

The $Contact Handler$ performs Bernoulli trials to decide whether a contact
between an infectious and susceptible person leads to disease transmission. 
People transit through Susceptible-Exposed-Infected-Recovered states,
similar to an influenza-like disease.
Each $ContactPool$ contains a link to its members and the $Population$ stores all personal
data, with $Person$ objects. We distinguish different a number of types of ContactPools associated with the household, the workplace, 
the school, \ldots. The Household is a key type because it fixed the home address of a person.
The household, workplace and school clusters are handled separately from the
community clusters, which are used to model general community contacts. The
$Population$ is a collection of $Person$ objects.



